%% ma: L12-B14-0001
%% lop: 12
%% chuong: 5
%% bai: 14
%% dang: 12.14.1
%% loai: TN4
%% muc: NB
%% dap_an: "B"
%% nguon: "0. ĐỀ THAM KHẢO TN THPT 2025.pdf (D00), Phần I Câu 7"
%% kien_thuc: "Bài 14 (vectơ pháp tuyến của mặt phẳng)"
%% trang_thai: chinh-thuc
\begin{ex}
Trong không gian với hệ trục tọa độ $Oxyz$, cho mặt phẳng $(P)$ có phương trình $x-3y-z+8=0$. Vectơ nào sau đây là một vectơ pháp tuyến của mặt phẳng $(P)$?
\choice{$\vec n_1(1;-3;1)$}{\True $\vec n_2(1;-3;-1)$}{$\vec n_3(1;-3;8)$}{$\vec n_4(1;3;8)$}
\loigiai{Mặt phẳng $ax+by+cz+d=0$ có một vectơ pháp tuyến $(a;b;c)$. Với $(P)$: $\vec n_2(1;-3;-1)$.}
\end{ex}
